% GENERATED FILE. Edit canonical lessons, then run npm run generate:books.
% canonical-source-hash: fnv1a64:6f7a8349cc0af6df
% canonical-lessons: TE-C291-daksinam, TE-C291-pradesam, TE-C291-sarihaddu, TE-C291-eduruga, TE-C291-cuttu2

\chapter{South, Place, Border, Opposite, Around}
\label{ch:te-south-place-border-opposite-around}
\hlchaptermodality{291}

\begin{chapteropening}
\textbf{By the end of this chapter:} \emph{I can say south, a place, a border, opposite and around.}

Complete the last lesson of chapter 291: I can say south, a place, a border, opposite and around.
\end{chapteropening}

\section[\texorpdfstring{\te{దక్షిణం} (dakṣiṇaṁ)}{దక్షిణం (dakṣiṇaṁ)}]{\te{దక్షిణం} (dakṣiṇaṁ) — south}

\label{lesson:TE-C291-daksinam}



\begin{quote}
Before the new one: say the Telugu for a century, then the Telugu for childhood.
\end{quote}

\subsection*{You'll want to know: \te{దక్షిణం}}
\textbf{\te{దక్షిణం}} — \emph{dakṣiṇaṁ} — \textquotedblleft{}south\textquotedblright{}.

\begin{grammarlens}[title={What you've built}]
One more word for this chapter.
\end{grammarlens}

\subsection*{Your turn}
\begin{itemize}
\raggedright
  \item \emph{Say it:} \emph{dakṣiṇaṁ}
  \item \emph{Say it:} \emph{dakṣiṇaṁ}, once more
  \item \emph{Say it:} say \emph{dakṣiṇaṁ}
  \item \emph{Recall:} say the Telugu for a century, then the Telugu for childhood, then say \emph{dakṣiṇaṁ} again
\end{itemize}

\begin{tcolorbox}[breakable,colback=teal!4,colframe=teal!35!black,title={Before you move on}]
What does \te{దక్షిణం} mean? (\textquotedblleft{}South\textquotedblright{}.) Say it once more.
\end{tcolorbox}
\section[\texorpdfstring{\te{ప్రదేశం} (pradēśaṁ)}{ప్రదేశం (pradēśaṁ)}]{\te{ప్రదేశం} (pradēśaṁ) — a place, a spot}

\label{lesson:TE-C291-pradesam}



\begin{quote}
Before the new one: say the Telugu for childhood, then the Telugu for south.
\end{quote}

\subsection*{You'll want to know: \te{ప్రదేశం}}
\textbf{\te{ప్రదేశం}} — \emph{pradēśaṁ} — \textquotedblleft{}a place, a spot\textquotedblright{}.

\begin{grammarlens}[title={What you've built}]
One more word for this chapter.
\end{grammarlens}

\subsection*{Your turn}
\begin{itemize}
\raggedright
  \item \emph{Say it:} \emph{pradēśaṁ}
  \item \emph{Say it:} \emph{pradēśaṁ}, once more
  \item \emph{Say it:} say \emph{pradēśaṁ}
  \item \emph{Recall:} say the Telugu for childhood, then the Telugu for south, then say \emph{pradēśaṁ} again
\end{itemize}

\begin{tcolorbox}[breakable,colback=teal!4,colframe=teal!35!black,title={Before you move on}]
What does \te{ప్రదేశం} mean? (\textquotedblleft{}A place, a spot\textquotedblright{}.) Say it once more.
\end{tcolorbox}
\section[\texorpdfstring{\te{సరిహద్దు} (sarihaddu)}{సరిహద్దు (sarihaddu)}]{\te{సరిహద్దు} (sarihaddu) — a border, a boundary}

\label{lesson:TE-C291-sarihaddu}



\begin{quote}
Before the new one: say the Telugu for south, then the Telugu for a place.
\end{quote}

\subsection*{You'll want to know: \te{సరిహద్దు}}
\textbf{\te{సరిహద్దు}} — \emph{sarihaddu} — \textquotedblleft{}a border, a boundary\textquotedblright{}.

\begin{grammarlens}[title={What you've built}]
One more word for this chapter.
\end{grammarlens}

\subsection*{Your turn}
\begin{itemize}
\raggedright
  \item \emph{Say it:} \emph{sarihaddu}
  \item \emph{Say it:} \emph{sarihaddu}, once more
  \item \emph{Say it:} say \emph{sarihaddu}
  \item \emph{Recall:} say the Telugu for south, then the Telugu for a place, then say \emph{sarihaddu} again
\end{itemize}

\begin{tcolorbox}[breakable,colback=teal!4,colframe=teal!35!black,title={Before you move on}]
What does \te{సరిహద్దు} mean? (\textquotedblleft{}A border, a boundary\textquotedblright{}.) Say it once more.
\end{tcolorbox}
\section[\texorpdfstring{\te{ఎదురుగా} (edurugā)}{ఎదురుగా (edurugā)}]{\te{ఎదురుగా} (edurugā) — opposite, facing}

\label{lesson:TE-C291-eduruga}



\begin{quote}
Before the new one: say the Telugu for a place, then the Telugu for a border.
\end{quote}

\subsection*{You'll want to know: \te{ఎదురుగా}}
\textbf{\te{ఎదురుగా}} — \emph{edurugā} — \textquotedblleft{}opposite, facing\textquotedblright{}.

\begin{grammarlens}[title={What you've built}]
One more word for this chapter.
\end{grammarlens}

\subsection*{Your turn}
\begin{itemize}
\raggedright
  \item \emph{Say it:} \emph{edurugā}
  \item \emph{Say it:} \emph{edurugā}, once more
  \item \emph{Say it:} say \emph{edurugā}
  \item \emph{Recall:} say the Telugu for a place, then the Telugu for a border, then say \emph{edurugā} again
\end{itemize}

\begin{tcolorbox}[breakable,colback=teal!4,colframe=teal!35!black,title={Before you move on}]
What does \te{ఎదురుగా} mean? (\textquotedblleft{}Opposite, facing\textquotedblright{}.) Say it once more.
\end{tcolorbox}
\section[\texorpdfstring{\te{చుట్టూ} (cuṭṭū)}{చుట్టూ (cuṭṭū)}]{\te{చుట్టూ} (cuṭṭū) — around}

\label{lesson:TE-C291-cuttu2}



\begin{quote}
Before the new one: say the Telugu for a border, then the Telugu for opposite.
\end{quote}

\subsection*{You'll want to know: \te{చుట్టూ}}
\textbf{\te{చుట్టూ}} — \emph{cuṭṭū} — \textquotedblleft{}around\textquotedblright{}.

\begin{grammarlens}[title={What you've built}]
One more word for this chapter.
\end{grammarlens}

\subsection*{Your turn}
\begin{itemize}
\raggedright
  \item \emph{Say it:} \emph{cuṭṭū}
  \item \emph{Say it:} \emph{cuṭṭū}, once more
  \item \emph{Say it:} say \emph{cuṭṭū}
  \item \emph{Recall:} say the Telugu for a border, then the Telugu for opposite, then say \emph{cuṭṭū} again
\end{itemize}

\begin{tcolorbox}[breakable,colback=teal!4,colframe=teal!35!black,title={Before you move on}]
What does \te{చుట్టూ} mean? (\textquotedblleft{}Around\textquotedblright{}.) Say it once more.
\end{tcolorbox}
